\section{Security Groups}

\slidebackground{backgrounds/wallpaper.jpg}

\begin{frame}[fragile]{What Is a Security Group?}
	\begin{itemize}
		\item A \textbf{security group} is a role stored in \texttt{res.groups}.
		\item Users are assigned to one or more groups.
		\item Groups control access rights, record rules, and UI visibility.
\begin{block}{Core Idea}
	Give permissions to groups, then assign users to groups -- not the opposite.
\end{block}
	\end{itemize}
\end{frame}

\begin{frame}[fragile]{Defining a Category and Groups}
\begin{lstlisting}[style=xml]
<record id="module_category_school" model="ir.module.category">
  <field name="name">School</field>
  <field name="sequence">20</field>
</record>

<record id="group_school_user" model="res.groups">
  <field name="name">User</field>
  <field name="category_id" ref="module_category_school"/>
  <field name="implied_ids" eval="[(4, ref('base.group_user'))]"/>
</record>

<record id="group_school_manager" model="res.groups">
  <field name="name">Manager</field>
  <field name="category_id" ref="module_category_school"/>
  <field name="implied_ids" eval="[(4, ref('group_school_user'))]"/>
  <field name="users" eval="[(4, ref('base.user_root')), (4, ref('base.user_admin'))]"/>
</record>
\end{lstlisting}
\end{frame}

\begin{frame}[fragile]{Breaking Down Group Fields}
	\begin{itemize}
		\item \textbf{\texttt{name}}: group label (User, Manager, Officer)
		\item \textbf{\texttt{category\_id}}: Settings app section for the group
		\item \textbf{\texttt{implied\_ids}}: groups automatically included
		\item \textbf{\texttt{users}}: optional default members
		\item \textbf{\texttt{comment}}: description of the role
	\end{itemize}
\begin{lstlisting}[style=xml]
<field name="implied_ids" eval="[(4, ref('group_school_user'))]"/>
\end{lstlisting}
	\begin{itemize}
		\item Manager implies User $\Rightarrow$ Manager gets all User rights too.
	\end{itemize}
\end{frame}

\begin{frame}[fragile]{Typical School Privilege Ladder}
	\begin{enumerate}
		\item \textbf{User}: read students, limited write
		\item \textbf{Officer/Teacher}: create/update operational data
		\item \textbf{Manager}: full CRUD + configuration
	\end{enumerate}
\begin{lstlisting}[style=xml]
<!-- Manager implies Officer implies User -->
<field name="implied_ids" eval="[(4, ref('group_school_officer'))]"/>
\end{lstlisting}
	\begin{itemize}
		\item Design roles as a ladder, not disconnected islands.
	\end{itemize}
\end{frame}

\begin{frame}[fragile]{Checking Groups in Python}
\begin{lstlisting}
if self.env.user.has_group('school_manager.group_school_manager'):
	# manager-only logic
	pass

if not self.env.user.has_group('school_manager.group_school_user'):
	raise AccessError("You need School User access.")
\end{lstlisting}
	\begin{itemize}
		\item Prefer XML-ID group checks over hard-coded group IDs.
		\item Still define real ACL/record rules; Python checks are extra safety.
	\end{itemize}
\end{frame}

\begin{frame}{Group Design Best Practices}
	\begin{itemize}
		\item Create a module category for your app.
		\item Keep group names short and clear.
		\item Use \texttt{implied\_ids} for privilege inheritance.
		\item Assign rights to groups, not individual users.
		\item Put group XML in \texttt{security/} and load it before access CSV.
	\end{itemize}
\end{frame}
